% All reported measurements are generated from committed CSVs by
% scripts/make_paper_numbers.py. Add a macro there; never type a result here.
\documentclass[11pt,a4paper]{article}
\PassOptionsToPackage{hyphens}{url}
\usepackage[margin=1in]{geometry}
\usepackage{amsmath,amssymb,array,booktabs,graphicx,xcolor,natbib}
\usepackage[hidelinks]{hyperref}
\usepackage[T1]{fontenc}
\usepackage{lmodern}
\usepackage{microtype}
\hypersetup{pdftitle={Credit Portfolio Optimization with cuOpt: A Reproducible Benchmark of Risk, Constraints, and Solver Endpoints},
 pdfauthor={Fred Hu},pdfsubject={Technical working paper, revision candidate},
 pdfkeywords={credit portfolios, cuOpt, CVaR, tracking error, feasibility, bond lots}}
\newcolumntype{L}[1]{>{\raggedright\arraybackslash}p{#1}}
\graphicspath{{../results/}{./}}
\bibliographystyle{plainnat}
\input{numbers}

\newcommand{\MaybeFigure}[2]{%
  \IfFileExists{../results/#1}{\includegraphics[width=\linewidth]{#1}}{%
  \IfFileExists{#1}{\includegraphics[width=\linewidth]{#1}}{%
  \fbox{\parbox{\dimexpr\linewidth-2\fboxsep-2\fboxrule\relax}{%
  \centering\small [figure \texttt{#1} unavailable: #2]}}}}}

\title{Credit Portfolio Optimization with cuOpt:\\
       A Reproducible Benchmark of Risk, Constraints, and Solver Endpoints}
\author{Fred Hu\thanks{Correspondence: \texttt{fred.qian.hu@gmail.com}. Artifacts:
\url{https://github.com/FredHuFuture/credit-cuopt-benchmark}.}}
\date{October 8, 2026\\\small Revision candidate; technical working paper}

\begin{document}
\maketitle

\begin{abstract}
We study NVIDIA cuOpt for continuous corporate-bond portfolio construction
under spread and default risk, rating and issuer limits, duration matching,
duration-times-spread bounds, and turnover constraints. The contribution is
an open application benchmark, not a new optimization algorithm. We compare
four endpoints: a raw first-order point, a feasibility-repaired point, a basic
solution, and a directly modeled conic solution. A central case is a
standard-deviation tracking-error target. A Gaussian-calibrated
mean-absolute-deviation (MAD) proxy set to the same nominal budget delivers
\PZeroNaiveMadSigmaBp\,bp against a ten-basis-point target, an overshoot of
\PZeroNaiveMadOvershoot$\times$. Recalibration is possible: a
\PZeroSelectedMadBudgetBp\,bp MAD budget delivers
\PZeroSelectedMadSigmaBp\,bp in-model and remains below target on
\PZeroValidationDrawCount{} fixed validation draws. It requires a
\PZeroSelectedMadHaircut$\times$ nominal haircut, reduces the scaled objective
by \PZeroSelectedMadObjectiveLossBp\,bp, and increases CVaR loss by
\PZeroSelectedMadCvarCostBp\,bp relative to the exact SOCP portfolio.
Thus the defect is not unavailability of the mandate but failure of naive
calibration, with an observable cost--accuracy trade-off.

Both conic solvers deliver the reference-size sigma mandate; working-size
Clarabel remains feasible while the native cuOpt path records a crash.
Contended SOCP timings are excluded. On the MAD-form LP, tight raw PDLP passes
the proxy's numerical constraints without crossover, but this is not sigma
compliance. A paired turnover leave-one-out and a matched-feasibility QP
comparison show that constraint structure, representation, and solver endpoint
matter alongside device choice. Results use a controlled credit scenario
generator, not an empirical asset-pricing test or an executable bond-trading
system. Supplemental pilots examine single-request timing and hypothetical
bond-lot repair, retaining failed quality gates rather than reporting new
speed advantages. Reported measurements are generated from archived result files.
\end{abstract}

\section{Introduction}

Scenario Mean-CVaR is attractive for credit because it represents asymmetric
spread widening and default losses directly, and because its finite-scenario
form is a linear program (LP)~\citep{rockafellar2000cvar}. The dense scenario
block is also a natural target for GPU first-order methods. That observation,
however, is only the beginning of a portfolio benchmark. Institutional
portfolios are judged on the final mandate, not on the objective value returned
by a solver at its native stopping tolerance.

Credit is not merely a new asset label on an equity example. Spread shocks
scale with spread level, multiple bonds share one issuer's default event, and
duration and DTS limits prevent a tail-risk optimizer from obtaining an
apparent improvement by abandoning the benchmark's risk exposures.
Scenario CVaR in credit is established~\citep{andersson2001credit}, as is
DTS-based risk measurement~\citep{bendor2007dts}. Our question is how a
cuOpt-based implementation behaves when these features are present together.

Three distinctions are easy to erase in a speed table. First, a sigma
tracking-error limit is a second-order cone (SOC), whereas its common MAD proxy
is an LP constraint. Second, an approximate first-order point, a repaired
continuous point, a basic solution, and an exact conic point are different
endpoints. Third, dense and factorized covariance representations change the
computational structure and memory footprint without changing the QP being
solved. This
paper makes those distinctions explicit and refuses a timing ratio whenever
the compared rows do not reach the same endpoint and quality target.

The contribution is an experimental application benchmark rather than a claim
to new cone geometry or new credit-risk theory. The
SOC form of tracking error is standard~\citep{jorion2003tracking}; MAD portfolio
models are classical~\citep{konno1991mad}; the ratio between absolute and
quadratic dispersion goes back to \citet{geary1935ratio}; and lifted
polyhedral approximations can enforce a cone conservatively with LP
machinery~\citep{bental2001soc,vielma2008lifted}. What is less often measured is
what happens after a portfolio optimizer is allowed to select into the proxy's
blind spot, how much recalibration costs, and whether solver speed survives an
independent mandate checker.

Our findings are fourfold.
\begin{enumerate}
\item Naive MAD calibration materially misses a sigma mandate. Recalibration
      restores delivery, but at a measurable CVaR and objective cost.
\item A tight raw first-order point satisfies the MAD-form LP without
      crossover. That proxy compliance must not be mistaken for delivery of
      the original sigma mandate.
\item Turnover is an important contributor to crossover cost, not its sole
      cause, according to a paired leave-one-out experiment.
\item Matched-feasibility QP results depend jointly on representation, solver
      stack, endpoint, and size. Rows repaired on CPU and raw on GPU are
      reported separately and never converted into a speed ratio.
\end{enumerate}

The main benchmark is a continuous relaxation. Bond lots, minimum trade sizes,
cardinality, transaction-cost schedules, and operational pre-trade checks are
distinct implementation requirements, as in institutional bond-optimization
work~\citep{benslimane2023bondmip}. Appendix~\ref{app:pilots} adds a small
hypothetical lot-repair pilot and a single-request timing protocol. Neither
establishes market execution readiness or extends the main performance grid.

\section{Model and mandate}

Let $w\in\mathbb{R}^N$ be portfolio weights, $w^0$ the incumbent, and
$R\in\mathbb{R}^{S\times N}$ scenario returns. With loss
$\ell_s(w)=-R_s w$, empirical CVaR at level $\alpha$ is represented by
\begin{equation}
 \operatorname{CVaR}_{\alpha}(w)=
 \min_{t,u}\left\{t+\frac{1}{(1-\alpha)S}\sum_s u_s:
 u_s\geq \ell_s(w)-t,\ u_s\geq0\right\}.
\end{equation}
The CVaR experiments maximize
\begin{equation}
 J(w)=\mu^\top w-\lambda\operatorname{CVaR}_{\alpha}(w),
 \qquad \alpha=0.95,\quad\lambda=3,
\end{equation}
subject to full investment, long-only name bounds, and the credit mandate.
$R$ and $\mu$ are monthly excess returns over a duration-matched Treasury
leg; $\mu$ is the generator's carry, convexity, and expected-default-loss
estimate, not a fitted alpha forecast. Turnover is a constraint,
$\frac12\lVert w-w^0\rVert_1\leq0.25$, not an objective penalty in these
experiments. Expected excess return and CVaR are reported in monthly basis
points; a ``scaled objective'' difference means $10^4\Delta J$, not an
observed return or a separate economic risk measure.

The credit limit sheet includes rating bands, sector and issuer
caps, cell and aggregate duration matching, duration-times-spread bounds, a
turnover budget, full investment, and name bounds. The artifact independently
recomputes every family from the returned weights. In the default mandate,
each rating bucket has a lower bound equal to its benchmark weight minus the
smaller of 0.10 and half that weight, and an upper bound 0.10 above benchmark
(clipped to one). Sector caps are benchmark plus 0.05. A common issuer cap is
the larger of 0.01 and 1.5 times the largest benchmark issuer weight; the name
cap is the larger of 0.005 and three times the largest benchmark name weight.
Aggregate duration stays within 0.10 years of benchmark, with 0.25-year
cell-duration bands, and DTS lies between 0.9 and 1.1 times benchmark DTS.
These are benchmark design choices, not universal desk standards.

For active weights $x=w-w^b$, let $A$ be the column-centered monthly scenario
matrix. The modeled annualized sigma tracking-error mandate is
\begin{equation}
  \sqrt{12/S}\,\lVert Ax\rVert_2\leq \tau,
  \qquad \Sigma_{\mathrm{ann}}=12A^\top A/S.
\end{equation}
This is an SOC constraint; the empirical variance uses divisor $S$, not
$S-1$. Budgets quoted in basis points are converted to return units before
optimization. The LP proxy replaces
quadratic dispersion by scaled sample MAD,
\begin{equation}
 \sqrt{12}\sqrt{\pi/2}\;S^{-1}\sum_s |(Ax)_s|\leq\tau.
\end{equation}
The Gaussian scaling is exact in population for a centered normal scalar;
it is not a finite-sample identity or a characterization of normality. It is neither
an inner nor an outer approximation for a general credit-return distribution.
Conservatively scaled polyhedral cone approximations provide another LP
route~\citep{bental2001soc}; we do not benchmark that alternative. Our MAD arm
tests a simple Gaussian-calibrated LP proxy, not a claim that desks universally
use this calibration.

\section{Experimental design}

\subsection{Data and scenario model}

The working problem has \MSixBonds{} synthetic bonds and
\MSixScenarios{} monthly scenarios, with universe seed 7 and scenario seed 11.
The main comparisons hold the universe, scenario seed, mandate, and objective
fixed. Scenarios combine correlated Student-$t$ spread factors, idiosyncratic
variation, asymmetric widening, and issuer-level jump-to-default. Spread
changes are proportional to initial spread; all bonds of an issuer share a
default event. Expected excess returns include spread carry and expected
default losses. The return approximation omits some pricing effects, including
the rate--spread convexity cross term. This is a controlled benchmark
generator, not an estimated production risk model.
Appendix~\ref{app:generator} specifies the operating parameters, benchmark,
incumbent, and maturity-cell constraints.

An optional SEC Form N-PORT pipeline supplies real holdings-derived security
characteristics, with spread-implied ratings and proxy sectors explicitly
labeled. It is an artifact extension, not the input to the main comparisons
and not evidence of real-market portfolio performance. Result filenames use
the internal prefix \texttt{p0\_} for the corrected experiment series.

The jump experiment is descriptive mechanism evidence. It connects naturally
to robust dispersion diagnostics~\citep{geary1935ratio} and to the distinction
between continuous variation and jumps in bipower variation
\citep{barndorff2004bipower}. We do not call the four-cell contrast a variance
decomposition or infer a structural causal share from it.

\subsection{Endpoints and independent gates}

Each row names a method class and endpoint:
\begin{center}
\begin{tabular}{lll}
\toprule
Endpoint & Construction & Interpretation\\
\midrule
raw & solver return & approximate continuous point\\
repaired & minimum-distance LP repair & continuous feasible point\\
basic & crossover or simplex & basic feasible point\\
exact-conic & SOCP barrier & sigma mandate represented directly\\
\bottomrule
\end{tabular}
\end{center}

The checker recomputes objective, expected excess return, CVaR, sigma TE, MAD
TE, bounds, and every mandate family from weights rather than trusting solver
metadata. A public timing ratio requires convergence, checker success, mandate
purity, objective consistency, and the same formulation and endpoint.
The purity gate allows at most $10^{-6}$ of each constraint's characteristic
scale, as defined in \path{src/credit.py}. The separate illustrative desk gate
allows 1bp of portfolio weight, 0.005 years of duration, 0.5 bp-years of DTS,
or 0.05bp annualized TE, depending on the quantity. It is not a claim of
industry acceptance. Legacy \texttt{respects\_mandate} fields in the TE
runner use a looser $10^{-3}$ relative threshold; publication decisions use
\texttt{purity\_pass} and the independently recomputed sigma target instead.
In a MAD row, purity checks the MAD constraint, not an additional sigma cap.
QP ratios match feasibility and objective-recomputation gates, not certified
identical distance to the mathematical optimum; objective gaps are disclosed
separately. DNF, native error, out-of-memory, and harness memory refusal are
recorded outcomes, not silently dropped cells.

\subsection{Timing discipline}

Main measurements were collected in August 2026 on an \EnvCpu{} CPU and an
\EnvGpu{} GPU (\EnvGpuMemory), under WSL2. The software stack uses Python
\EnvPython, CVXPY \EnvCvxpy, cuOpt \EnvCuopt, Clarabel 0.11.1, and HiGHS
\EnvHighs; the archived WSL environment reports 6.8GiB of visible RAM.
\path{requirements.lock} preserves the measured package versions.

Build, first solve, steady-state median, repair, and total endpoint time are
separate fields. Runs capture GPU utilization, memory, and process context. The
reference-size SOCP file explicitly flags GPU contention. The working-size
Clarabel SOCP repeats also ran with foreground workload and swap pressure.
Both files support feasibility only: no SOCP speed ratio is published.
The earlier draft's reference-size ratio has been withdrawn. Other reported
timings have clean file-level capture flags, which are snapshots rather than
proof of continuous system idleness. Runs generally comprise three repeats;
repair and leave-one-out use five. Failed cells are retained. The TE runner
stops after its first failed attempt, but the archived endpoint-ladder HiGHS
simplex arm includes three timed capped runs; the no-repeat policy was not
uniformly implemented. The artifact distinguishes
within-process steady-state repeats from process-isolated experiments; these
timings are descriptive workstation measurements, not confidence intervals.
The October supplemental pilot uses a separate single-request clock and stores
its own environment and source hashes. Its observations are not pooled with
the August timings; see Appendix~\ref{app:request}.

\section{Exact SOCP and MAD recalibration}

\subsection{Native conic validation}

At reference size, both conic stacks deliver the target independently:
\begin{center}
\begin{tabular}{lr}
\toprule
Stack & achieved sigma TE (bp)\\
\midrule
cuOpt barrier & \PZeroSocpRefCuoptTeBp\\
Clarabel & \PZeroSocpRefClarabelTeBp\\
\bottomrule
\end{tabular}
\end{center}
This validates native cuOpt conic feasibility on the 300-bond,
3,000-scenario reference problem, not relative speed. ``Exact'' here means
the sigma constraint is represented directly; the solver still uses numerical
tolerances.

At working size, Clarabel still returns a portfolio with
\PZeroSocpWorkingClarabelTeBp\,bp sigma TE and expected excess return of
\PZeroSocpWorkingClarabelExpectedExcessBp\,bp per month. The cuOpt barrier path
records \texttt{\PZeroSocpWorkingCuoptStatus} in the native sparse factorization
path. We ran that failure once, retained the status, and did not convert it into
a fictitious lower-bound timing. Conversely, we exclude the contaminated
Clarabel repeats from speed claims. The working-size evidence therefore answers
feasibility, not comparative timing.

\subsection{Naive versus recalibrated MAD}

\begin{center}
\small
\begin{tabular}{lrrL{3.0cm}}
\toprule
Method & budget (bp) & sigma TE (bp) & interpretation\\
\midrule
naive MAD & \MSixMadTenBudgetBp & \PZeroNaiveMadSigmaBp &
\PZeroNaiveMadOvershoot$\times$ target\\
recalibrated MAD & \PZeroSelectedMadBudgetBp & \PZeroSelectedMadSigmaBp &
passes true sigma gate\\
exact SOCP & \MSixMadTenBudgetBp & \PZeroSocpWorkingClarabelTeBp &
mandate represented directly\\
\bottomrule
\end{tabular}
\end{center}

The selected MAD budget is not the largest nominal budget. Among candidates
that pass sigma TE and every mandate gate, it is the one with the highest
independently recomputed objective. Its sigma TE ranges from
\PZeroValidationSigmaMinBp\,bp to \PZeroValidationSigmaMaxBp\,bp on
\PZeroValidationDrawCount{} fixed independent scenario draws, none used to
choose the budget.

Candidate generation begins at MAD budgets 1, 2, 3, 5, and 10bp, testing
PDLP tolerances $10^{-4}$ and $10^{-6}$. Adjacent achieved-sigma values bracket
the 10bp target between 2 and 3bp. Bisection adds midpoint budgets, up to eight
steps or an interval width of 0.1bp; selection uses all eligible recorded
candidates rather than assuming their objective is monotone. Validation
seeds 101, 211, and 307 score the fixed selected weights on fresh 20,000-row
draws and do not enter selection. This is not the best possible recalibrated
MAD frontier: a tighter solve of a rejected candidate could improve the cost.

The recalibration is therefore feasible but not free. Relative to exact SOCP,
expected excess return decreases by
\PZeroSelectedMadExpectedExcessLossBp\,bp per month, CVaR loss increases by
\PZeroSelectedMadCvarCostBp\,bp, and the scaled objective decreases by
\PZeroSelectedMadObjectiveLossBp\,bp. The defensible
conclusion is a cost--accuracy trade-off: naive calibration fails, conservative
recalibration works on these draws, and exact conic modeling avoids having to
learn a proxy haircut.

\subsection{Why defaults matter}

The separate tail sweep uses 120 bonds, 4,000 scenarios, universe seed 7,
and scenario seeds 11, 12, and 13. It keeps the optimization and mandate fixed
while varying continuous-tail degrees of freedom and default intensity.
At the 5bp MAD budget, defaults-off per-degree-of-freedom median overshoot
ranges from \MNineNoJumpOvershootMin$\times$ to
\MNineNoJumpOvershootMax$\times$; the minimum defaults-on cell median is
\MNineDefaultsOnOvershootMin$\times$. This is not a lower bound on individual
draws or other budgets. At $\nu=5$, raising the default-probability multiplier
from 0.01 to 0.30 raises median overshoot from
\MNineDoseOneOvershoot$\times$ to \MNineDoseThirtyOvershoot$\times$.
This supports the mechanism that an
absolute-deviation proxy underweights rare large losses relative to quadratic
dispersion. It does not identify a universal causal percentage.

\begin{figure}[tp]
\centering
\MaybeFigure{m9_tail_sweep.pdf}{run \texttt{scripts/plot\_tail\_sweep.py}}
\caption{Paired tail and jump controls. Panel A shows cell medians and
min--max ranges across scenario draws. Panel B distinguishes individual draws
from cell medians; panel C compares optimized and fixed-book Gaussian
controls across scenario counts. The experiment is mechanism evidence, not
a structural decomposition.}
\label{fig:tail}
\end{figure}
\clearpage

\section{Endpoint ladder}

The endpoint ladder asks whether satisfying the \emph{MAD-form LP} requires
a basic solution. The table uses the nominal MAD@10bp problem, not the
recalibrated proxy and not an exact sigma constraint.
\begin{center}
\small
\begin{tabular}{L{3.0cm}rrL{3.3cm}}
\toprule
Endpoint & reported time (s) & gate verdict & interpretation\\
\midrule
PDLP raw, loose & \PZeroEndpointLooseRawSeconds &
\PZeroEndpointLooseRawVerdict & approximate point\\
loose + L1 repair & \PZeroEndpointLooseRepairSeconds &
\PZeroEndpointLooseRepairVerdict & repaired continuous point\\
PDLP raw, tight & \PZeroEndpointTightRawSeconds &
\PZeroEndpointTightRawVerdict & MAD-form feasible point\\
crossover/basic & $>\PZeroEndpointBasicLowerBoundSeconds$ &
\PZeroEndpointBasicStatus & no completed endpoint\\
\bottomrule
\end{tabular}
\end{center}

Successful rows report steady-state solve medians, plus a separately measured
repair median where applicable. They exclude scenario generation, model
construction, and the first solve. The crossover row instead reports elapsed
time to timeout without a completed endpoint, not a steady-state median.
These are not end-to-end cold-start latencies.

The loose repair clears the desk gate but not the stricter purity gate because
the TE proxy itself is not visible to a repair over the linear mandate
families. Tightening PDLP produces a raw point that passes both LP gates.
Its true sigma TE remains \PZeroNaiveMadSigmaBp\,bp, so it \emph{fails} the
10bp sigma target. The conclusion is solely that the MAD-form LP can be
satisfied by a continuous nonbasic point; crossover is not necessary for
that narrower endpoint.
A basic solution may still be useful for downstream operations, but this
continuous endpoint ladder does not model lots or claim that a vertex is a
desk requirement. The separate lot pilot in Appendix~\ref{app:lots} checks
quantities explicitly rather than treating a basis as a proxy for them.

\section{What contributes to crossover cost?}

The earlier unpaired ablation could not support a causal attribution to
turnover. We rerun the full linear credit mandate (without a TE constraint)
and a leave-one-out mandate in the same
session, with all other families unchanged. Crossover overhead is the
difference between crossover-enabled and raw solve medians, not an internal
solver phase timer.
\begin{center}
\begin{tabular}{lrr}
\toprule
Constraint set & auxiliary variables added & crossover overhead (s)\\
\midrule
budget and box & -- & \PZeroLZeroCrossoverSeconds\\
full mandate & \PZeroFullAuxVariables & \PZeroFullCrossoverSeconds\\
full minus turnover & \PZeroNoTurnoverAuxVariables &
\PZeroNoTurnoverCrossoverSeconds\\
\bottomrule
\end{tabular}
\end{center}

Removing turnover reduces overhead by \PZeroTurnoverMarginalSeconds\,s,
\PZeroTurnoverMarginalSharePct{} of the full value. Under the pre-specified
rule the verdict is \texttt{\PZeroTurnoverVerdict}: turnover contributes
materially, while the remaining constraints still leave crossover far above
the budget-and-box baseline. The percentage is a paired marginal contrast on
this instance, not a portable causal share.

\section{Matched-feasibility QP benchmark}

The companion QP branch maximizes $\mu^\top w-(\lambda/2)w^\top\Sigma w$
under the same non-TE credit limit sheet. Its factor-plus-diagonal covariance
is estimated using the scenario generator's factor structure, not asserted
equal to the full empirical covariance used by the sigma-TE experiment. A dense
$N\times N$ covariance and its factorized representation are mathematically
equivalent but computationally different. This distinction predates GPU
portfolio solvers~\citep{perold1984large} and remains central in current bond
work~\citep{benslimane2026bond}.

We compare strict Clarabel CPU and cuOpt GPU barrier stacks. Clarabel uses
$10^{-9}$ feasibility and absolute/relative gap tolerances with up to 2,000
iterations; cuOpt is pinned to barrier at its solver defaults. Raw rows require
the $10^{-6}$ mandate gate and relative solver-versus-recomputed objective
agreement within $10^{-6}$. Clarabel rows that miss the independent mandate
gate receive a minimum-distance repair. Repaired weights are checked for
feasibility and their QP objective is recomputed, but agreement with a
solver-reported objective for the original QP is not applicable at that new
point. Its zero discrepancy field is a self-comparison, not independent
solver confirmation; raw diagnostics remain archived. Repaired common-quality
time is the solve median plus repair time. Because a repaired CPU point
and raw GPU point are different endpoints, those pairs receive no ratio. Of
\PZeroQpEligibleRows{} independently eligible rows, only
\PZeroQpRatioRows{} form matched raw-endpoint ratios.
The table divides GPU by CPU steady-state solve medians; it excludes build,
scenario generation, and first-solve costs.

\begin{center}
\begin{tabular}{rrrr}
\toprule
Names & dense GPU/CPU & factored GPU/CPU & endpoint\\
\midrule
\PZeroQpFiveHundredBonds & \PZeroQpRawRatioDenseFiveHundred &
\PZeroQpRawRatioFactoredFiveHundred & raw/raw\\
\PZeroQpTwoThousandBonds & \PZeroQpRawRatioDenseTwoThousand &
\PZeroQpRawRatioFactoredTwoThousand & raw/raw\\
\PZeroQpThreeThousandBonds & \PZeroQpRawRatioDenseThreeThousand &
\PZeroQpRawRatioFactoredThreeThousand & raw/raw\\
\bottomrule
\end{tabular}
\end{center}

The signs differ by representation and size: there is no single ``GPU
speedup.'' The maximum relative gap to the best feasible objective among
eligible rows is \PZeroQpMaxObjectiveGap, which is reported rather than hidden
behind feasibility. Factorized GPU rows solve through
\PZeroQpLargestSolvedBonds{} names. At the \PZeroQpAmundiBonds{}-name case, dense
arms are refused by the harness's host-memory policy and factored arms record
\texttt{\PZeroQpAmundiFactoredStatus} during scenario generation. Dense does
not cease to exist mathematically; these are explicit resource outcomes of
this workstation and harness.

The comparison is adjacent to GPU conic work such as
CuClarabel~\citep{chen2024cuclarabel}, GPU portfolio solvers such as
FlashFolio~\citep{jiang2026flashfolio}, and deadline-oriented systems such as
PRISM~\citep{ghosh2026prism}. It differs by making endpoint and independent
mandate delivery prerequisites for a ratio.

\section{Threats to validity}

\paragraph{Model risk.} The scenario generator is parameterized using stylized
assumptions, not fitted to a market return panel or validated as a production
fixed-income risk model. N-PORT supplies real instruments and
characteristics, not transaction histories or an authoritative covariance.
The three validation draws test generator robustness, not market validity.

\paragraph{External validity.} Results come from one consumer GPU, one CPU,
one software stack, and a limited grid. Solver-stack ratios are not hardware
causal effects. A native crash is a version-specific recorded outcome, not a
claim that the cuOpt product class cannot represent SOCP.

\paragraph{Endpoint scope.} The main experiments study continuous portfolios.
Basic solutions, lots, cardinality, and executable trade lists are distinct
requirements. The endpoint ladder answers MAD-form LP feasibility only, not
sigma delivery. The supplemental MILP repairs linear and hypothetical quantity
constraints, not a sigma cone. A feasible capped incumbent is not certified
optimal. Neither pilot includes real security-level trading conventions,
liquidity or execution costs.

\paragraph{Proxy scope.} Recalibration is selected in-model and checked on
fixed independent draws. It may drift across regimes. An exact SOCP or a
conservative lifted polyhedral approximation avoids assuming a stable MAD to
sigma conversion.

\paragraph{Data rights.} The committed real-instrument universe is derived
from public-domain SEC N-PORT data. Licensed TRACE, WRDS, FISD, ICE, and ETF
holdings extracts are not redistributed. Optional runtime downloads retain
their source terms and hashes.

\section{Related work}

Tracking-error-constrained portfolio geometry is established by
\citet{jorion2003tracking}, while MAD optimization follows
\citet{konno1991mad}. The present contribution is not either formulation but
their measured mandate delivery after optimization and the resulting
cost--accuracy comparison. Geary's ratio~\citep{geary1935ratio} and bipower
variation~\citep{barndorff2004bipower} supply the statistical ancestry for the
absolute-versus-quadratic dispersion gap.

SOC constraints need not force a choice between an exact conic solver and a
nonconservative MAD proxy. Lifted polyhedral approximations are classical
\citep{bental2001soc,vielma2008lifted}; they are a natural extension because
this benchmark separately measures the cost of auxiliary-variable
reformulations.

NVIDIA's cuFOLIO provides general cuOpt portfolio workflows
\citep{nvidia2026cufolio}. CuClarabel accelerates conic
optimization on GPUs~\citep{chen2024cuclarabel}; FlashFolio studies a GPU
portfolio solver~\citep{jiang2026flashfolio}; PRISM studies deadline-bounded
CPU/GPU rebalancing~\citep{ghosh2026prism}; and \citet{niu2026meanvariance}
analyze large mean-variance problems. Amundi's bond work emphasizes both
security-level constraints, L1/L2 risk, and formulation choice
\citep{benslimane2026bond,benslimane2023bondmip}. Those results mean that
neither bond optimization nor the L1/L2 distinction is new. Our addition is
an open cuOpt-based credit benchmark joining issuer-level default scenarios,
credit exposure constraints, proxy recalibration, and solver-endpoint
measurements. These adjacent works do not establish that this exact
combination has already been benchmarked; equally, we do not claim an
exhaustively verified ``first'' implementation. The contribution is the
credit-specific experimental artifact and its findings, not the use of a
GPU or a feasibility checker alone.

\section{Conclusion}

cuOpt can solve useful continuous credit-portfolio formulations, but the answer
depends on risk modeling, constraints, and the required endpoint.
Naive MAD is fast but does not deliver the nominal sigma mandate.
Recalibration delivers it on the tested draws, with a
\PZeroSelectedMadHaircut$\times$ haircut and measurable economic cost. Exact
SOCP delivers the modeled mandate directly at reference size in both stacks;
no comparative SOCP speed claim survives the timing-contamination gate.
The working-size cuOpt native path records a crash and Clarabel supplies
feasibility evidence. Tight raw PDLP can satisfy the MAD-form LP without
crossover, but still violates the original sigma target. Removing turnover
reduces paired crossover overhead by \PZeroTurnoverMarginalSharePct;
substantial overhead remains.
On the QP branch, representation and endpoint matching prevent most tempting
CPU/GPU ratios from being publishable.

The practical recommendation is simple: define the required endpoint, score it
independently, and only then time it. Hardware accelerates a formulation; it
does not certify that the formulation delivered the mandate.

\section{Reproducibility}

The repository commits result CSVs, provenance headers, scenario seeds,
independent check fields, and recorded non-results. The paper never contains a
typed empirical result: \path{scripts/make_paper_numbers.py} generates
\path{paper/numbers.tex} and exits nonzero if a required cell is absent.
The principal evidence files are:
\begin{itemize}
\item \path{p0_socp_reference.csv} and \path{p0_socp_working.csv};
\item \path{p0_te_recalibration.csv};
\item \path{p0_endpoint_ladder.csv};
\item \path{p0_ablation_loo.csv};
\item \path{p0_qp_matched.csv}.
\item \path{m9_tail_sweep.csv} for the separate mechanism experiment.
\item \path{p1_bond_lots_cpu.csv}, with metadata and portfolio sidecars;
\item \path{p1_request_timing_cpu.csv} and
      \path{p1_request_timing_gpu_smoke.csv} for the supplemental pilots.
\end{itemize}
The schemas distinguish solver-returned outcomes from author classifications
where recorded; fields and completeness vary across runners. Historical
results remain in the repository but do not override the corrected
measurements. The release README gives separate CPU-only artifact checks
and optional GPU experiment reproduction. Original measurements are retained;
reruns should write to new paths. A passing number-generation step establishes
traceability, not by itself scientific validity or freedom from contention.

\appendix
\section{Synthetic model specification}\label{app:generator}

The operating parameters below are stylized assumptions in
\path{src/scenarios.py}, not estimates fitted to a bond-return panel.
Systematic level, slope, and credit shocks are standardized multivariate
Student-$t$ with $\nu=5$. Positive credit shocks are multiplied by 1.35, then
recentered; the widening transform therefore changes realized volatility.
Rating, sector, and bond residual shocks are Gaussian and demeaned in-sample.
Before the high-yield multiplier, their monthly relative-spread scales are
0.04, 0.05, and 0.06, and the systematic credit scale is 0.08. BB and B bonds
receive a 1.20 multiplier. Rate level and slope scales are 25 and 12bp;
their correlation is $-0.30$, and level--credit correlation is $-0.35$.
The reported excess-return branch cancels the first-order Treasury rate leg.

For spread $s_i$, spread duration $D_i$, convexity $C_i$, and monthly
absolute spread change $\Delta s_{si}$, the nondefault return approximation is
\[
 R_{si}=s_i/12-D_i\Delta s_{si}+\tfrac12 C_i(\Delta s_{si})^2.
\]
Default replaces this return with $(40-P_i)/P_i$, where $P_i$ is price per
100 face. Issuer asset shocks combine the standardized, sign-reversed credit
factor and an independent normal shock with loading correlation 0.15.
Thresholds are estimated from a separate Monte Carlo draw under that same
non-Gaussian distribution. Annual default probabilities for AAA, AA, A, BBB,
BB, and B are respectively 0.01\%, 0.03\%, 0.08\%, 0.25\%, 1.10\%, and
4.20\%; monthly targets divide these by 12. Every bond of an issuer shares
the same event. The drift $\mu_i$ is carry plus the sample second-order
convexity gain plus target default probability times default return, not a
sample-mean alpha estimate. Full details and diagnostic disagreements between
this drift and the realized sample mean are retained in the artifact.

\path{src/universe.py} generates fixed-rate bullet bonds. Benchmark weights
are proportional to issue size times price. Issue sizes follow a lognormal
distribution with median USD600 million, log standard deviation 0.60, and a
USD250 million floor. The incumbent samples 65\% of the bonds without
replacement using benchmark-weight probabilities; held benchmark weights are
multiplied by lognormal tilts with log standard deviation 0.20 and normalized.
Issuer rating probabilities for AAA through B are 1\%, 7\%, 32\%, 45\%,
10\%, and 5\%; rating OAS anchors are 45, 60, 90, 140, 260, and 420bp before
sector, issuer, tenor, and bond multipliers. These are design inputs, not
observations of a particular index.

Maturity buckets are $[0,3)$, $[3,7)$, $[7,12)$, and $[12,40)$ years.
For bucket $c$, its duration constraint is
\[
 \left|\sum_{i\in c}D_i(w_i-w_i^b)\right|\leq0.25\text{ years}.
\]
This bounds a contribution to total portfolio duration, not the bucket's
average duration. The benchmark and incumbent, unlike solver choice, remain
fixed within each paired experiment.

\section{Supplemental implementation pilots}\label{app:pilots}

These October 2026 pilots extend the artifact, not the main performance grid.
They use synthetic instruments and the same scenario generator, with universe
seed 7 and scenario seed 11. They are functional measurements on small
instances, not empirical validation, a scaling study or additional speedup
claims. Inspired by quantity-aware bond optimization
\citep{benslimane2023bondmip} and explicit timing/quality comparisons
\citep{ghosh2026prism,chen2024cuclarabel}, they do not reproduce those systems,
use a licensed risk model, or run CuClarabel.

\subsection{One request, one timing boundary}\label{app:request}

The new runner records build, explicit CVXPY canonicalization, solve call,
repair (zero in this pilot), and independent validation separately. A single
wall clock covers the whole request, including interstage overhead and
resource snapshots. Scenario generation, worker startup and checkpoint I/O
are excluded and identified separately. The solve-call bucket includes adapter,
transfer and result unpacking; it is not GPU kernel time. This differs from
legacy batch fields that can include a first solve plus multiple repeats.

Each solver starts in its own process. Its first request builds a fresh model;
the subsequent cached request, if accepted, reuses that CVXPY Problem and
requests warm start without asserting solver-state reuse. Cold means fresh
model, not necessarily cold libraries or device. A process watchdog caps each
request at 30s. All outcomes checkpoint immediately into a new output file;
failed acceptance suppresses cached requests rather than repeating failures.

The instance has 100 bonds and 1,000 CPU-generated scenarios, a 10bp annualized
sigma target, $\alpha=0.95$ and $\lambda=3$. The LP stacks solve MAD@10bp;
conic stacks model sigma directly. In addition to the original mandate purity
gate, acceptance requires budget/bound checks, a converged status, absolute
solver-versus-recomputed objective discrepancy at most $10^{-6}$, and true
sigma no more than $10+10^{-5}$bp. Clarabel uses $10^{-10}$ feasibility and gap
tolerances and one thread; HiGHS simplex uses $10^{-9}$ primal/dual feasibility
tolerances and one thread. PDLP uses the named tolerance and cuOpt barrier
retains its solver defaults. The archived configurations and package versions
are authoritative; the pilot does not silently equalize solver settings.

\begin{table}[htbp]
\centering\small
\begin{tabular}{L{4.1cm}rrL{3.4cm}}
\toprule
Request & wall (s) & sigma TE (bp) & acceptance\\
\midrule
HiGHS MAD, cold & \POneCpuHighsSeconds & \POneCpuHighsSigmaBp & rejected: true sigma\\
Clarabel SOC, cold & \POneCpuSocpColdSeconds & \POneCpuSocpColdSigmaBp & accepted\\
Clarabel SOC, cached & \POneCpuSocpCachedSeconds & \POneCpuSocpCachedSigmaBp & accepted\\
PDLP $10^{-4}$ MAD, cold & \POneGpuLooseSeconds & \POneGpuLooseSigmaBp & rejected: purity, sigma\\
PDLP $10^{-6}$ MAD, cold & \POneGpuTightSeconds & \POneGpuTightSigmaBp & rejected: purity, sigma\\
cuOpt SOC, cold & \POneGpuSocpSeconds & \POneGpuSocpSigmaBp & rejected: objective\\
\bottomrule
\end{tabular}
\caption{Single-request diagnostics, not comparable speed measurements. Each
entry is one observation, not a median. Different formulations and request
modes, and failed acceptance gates, preclude a comparative ratio. All solver
statuses are \texttt{optimal}; independent acceptance is a separate test.}
\label{tab:request}
\end{table}

The CPU file retains \POneCpuRowCount{} rows and the GPU file
\POneGpuRowCount{} rows. Both PDLP points pass the illustrative desk gate but
miss strict MAD purity and, separately, the true sigma target. The cuOpt SOC
portfolio satisfies budget, bounds, purity and sigma. Its objective discrepancy
is $\POneGpuSocpObjectiveGap$, just above the $10^{-6}$ gate. Its independently
recomputed objective differs from the cold same-model Clarabel reference by
only $\POneGpuSocpVsCpuObjectiveGap$; that closeness does not override the failed
gate or certify optimality. The CSV preserves weights but not the CVaR
auxiliary values. Default tolerance and auxiliary residuals are hypotheses for
further diagnosis, not established causes. No GPU row is accepted
(\POneGpuAcceptedCount{} accepted, \POneGpuRatioCount{} reported ratios).

The GPU run began after consecutive device-utilization samples passed a
25\% launch guard. This is not continuous workload exclusion: the per-request
telemetry consists only of device-wide before/after snapshots, not peaks or
proof of isolation. The table therefore provides a reproducible timing
\emph{boundary} and quality audit, not a new CPU/GPU or warm-start speed claim.

\subsection{Cash-aware lot repair}\label{app:lots}

The CPU quantity pilot uses 60 bonds, 500 scenarios and NAVs of USD1, 10 and
100 million. Hypothetical conventions are a USD1,000 par increment, USD10,000
minimum nonzero holding, USD5,000 minimum nonzero trade and a 2\% cash cap.
They are design inputs, not asserted market conventions. Price is a settlement
proxy with zero accrued interest. For integer lot counts $k_i$, par increment
$L$, price $P_i$ per 100 par, and NAV $V$,
\[
 w_i=k_iLP_i/(100V),\qquad \sum_iw_i+c=1,\qquad
 \tfrac12(\lVert w-w^0\rVert_1+|c-c^0|)\leq0.25.
\]
Cash $c$ is not renormalized away and earns zero excess return. The incumbent
is floored onto the lot lattice, sub-minimum positions removed and residual
cash retained. Its adjustment changes with NAV, so this is not a causal test
of lot size alone. The benchmark remains unchanged.

Within each NAV case the continuous cash-aware SOCP, naive rounding and repair
share the same mandate and starting book. HiGHS MILP minimizes bond-plus-cash
L1 distance from the continuous source subject to the linear mandate and
quantity rules. Separate buy/sell binaries prevent opposing trades from
evading minimum sizes. Both minimum sizes must be integer multiples of the
increment in this implementation. Crucially, this repair does not enforce the
sigma cone; true sigma is rescored afterward. Each solve has a 30s cap.

\begin{table}[htbp]
\centering\small
\begin{tabular}{lrrL{5.5cm}}
\toprule
NAV & rounded sigma & repaired sigma & repair outcome\\
\midrule
USD1M & \POneLotOneRoundedSigmaBp & \POneLotOneRepairSigmaBp & L1 optimal; sigma fails\\
USD10M & \POneLotTenRoundedSigmaBp & \POneLotTenRepairSigmaBp & capped feasible incumbent; not certified optimal\\
USD100M & \POneLotHundredRoundedSigmaBp & \POneLotHundredRepairSigmaBp & L1 optimal; strict sigma fails\\
\bottomrule
\end{tabular}
\caption{Quantity pilot. Sigma is annualized bp against a 10bp target.
All MILP repair rows pass the independent linear and quantity gates.
The complete artifact retains \POneLotRowCount{} endpoint rows, including the
three continuous references and all infeasible rounded portfolios.}
\label{tab:lots}
\end{table}

Rounding is not feasible in any case. At USD10M it passes sigma but overspends
cash and creates a sub-minimum trade. The repaired USD10M incumbent passes
all independent portfolio checks, but the solve reaches its time cap with a
\POneLotTenRepairGapPct{} MIP gap on L1 distance, not on economic objective.
It is not accepted as a converged endpoint. The USD100M repair misses sigma
by only \POneLotHundredRepairSigmaExcessBp\,bp, below the illustrative desk
tolerance but above the unchanged strict numerical gate. Solver optimal status
on the other MILP rows refers to the L1 repair within configured MIP tolerances,
not the original CVaR objective or executable trading quality. No lot row enters
a speed ratio. The pilot shows why linear/quantity compliance does not imply
delivery of a tight sigma budget; market-calibrated risk, security-specific
conventions, transaction costs, cardinality and MISOCP remain outside scope.

{\small\bibliography{references}}
\end{document}
